\section{Documentación, conclusiones y recomendaciones}

\subsection{Conclusiones}

\begin{pendiente}
Completar con las conclusiones derivadas de los resultados obtenidos. Deben
responder explícitamente a la pregunta de decisión formulada en la
sección~1.2 y vincularse con los tres desenlaces previstos.
\end{pendiente}

\subsection{Recomendaciones}

\begin{pendiente}
Formular las recomendaciones dirigidas al servicio, distinguiendo entre
aquellas que el modelo respalda con evidencia estadística y aquellas que
constituyen sugerencias de prudencia.
\end{pendiente}

\subsection{Limitaciones del estudio}

\begin{itemize}
  \item \textbf{Origen de los datos.} Con excepción de la tasa de arribos,
  consultada con una profesional del servicio, los valores numéricos son
  estimaciones del grupo y no provienen de registros del hospital. Los
  tiempos de servicio, las probabilidades de internación y las tasas de error
  de clasificación constituyen supuestos.
  \item \textbf{Horizonte de simulación.} La licencia educativa limita las
  corridas con Pedestrian Library a cinco horas de tiempo simulado, lo que
  obligó a introducir el factor de escala descripto en la
  sección~\ref{sec:escala}.
  \item \textbf{Velocidad de desplazamiento.} El escalado temporal exige
  multiplicar la velocidad de los peatones por el mismo factor, lo que
  distorsiona levemente la dinámica de los cuellos de botella físicos. Dado
  que los desplazamientos representan segundos frente a esperas de decenas de
  minutos, el efecto sobre las métricas se considera despreciable.
  \item \textbf{Alcance.} La exclusión de traumatología, radiografías y
  medicación concentra en los consultorios generales una demanda que el
  sistema real reparte.
  \item \textbf{Ausencia de abandonos.} El modelo no representa a los
  pacientes que se retiran sin ser atendidos.
  \item \textbf{Régimen estacionario.} Con utilizaciones cercanas a la unidad
  y un horizonte acotado, el sistema puede no alcanzar régimen estacionario,
  de modo que los promedios dependen de la duración de la corrida.
\end{itemize}

\subsection{Trabajo futuro}

\begin{itemize}
  \item Incorporar un segundo escenario alternativo consistente en la
  derivación de pacientes verdes a centros de atención primaria.
  \item Hacer depender la probabilidad de error de clasificación de la carga
  de trabajo del puesto de triage, de modo que el escenario alternativo
  incida también sobre la calidad de la clasificación y no sólo sobre los
  tiempos.
  \item Sustituir la tasa de arribos constante por una tasa variable por
  franja horaria.
  \item Incorporar el abandono de pacientes en espera.
  \item Validar los tiempos de servicio y las probabilidades de internación
  contra registros del hospital.
\end{itemize}

\subsection{Datos que permitirían refinar el modelo}

En caso de obtenerse acceso a registros del servicio, los datos de mayor
valor para el estudio serían los siguientes, ordenados por impacto:

\begin{enumerate}
  \item Marcas de tiempo de arribo de cada paciente, sobre varias semanas, que
  permitirían verificar el supuesto de tasa constante y estimar la variación
  horaria.
  \item Duración registrada del triage y de la consulta médica, discriminada
  por color.
  \item Proporción de pacientes derivados a traumatología, radiografías y
  medicación, que determinaría el caudal efectivo hacia los consultorios
  generales.
  \item Destino final por color, que permitiría cerrar la discusión sobre la
  probabilidad de internación del código naranja.
  \item Proporción de pacientes que se retiran sin ser atendidos.
\end{enumerate}

Debe señalarse que ninguno de estos datos requiere identificación de
pacientes: marcas de tiempo, color, consultorio y destino son suficientes, lo
que facilita su obtención desde el punto de vista del resguardo de la
información personal.

